\section{Methods}
\label{sec:methods}


\subsection{Data Collection}
\label{sec:methods-data-collection}
\subsubsection{ASB 018 standard}
To understand ASB 018's goals, we drew on both the standard itself and the factsheet that the ASB provides to accompany the standard.

\subsubsection{PGS audits (i.e., internal validation studies)}
To identify existing audit practices, we gathered PGS audit reports from the American Society of Crime Laboratory Directors Validation and Evaluation Repository\footnote{The repository can be accessed at this link: \url{https://www.ascld.org/validation-evaluation-repository/}.}. This database is, to our knowledge, the only publicly available centralized and voluntary repository for U.S. forensic labs' PGS audits. From this repository, we found five PGS audit reports, all for the STRmix PGS: the NYC Office of the Chief Medical Examiner (OCME) Lab's internal validation of STRmix v2.4~\cite{ocmeIV} and STRmix v2.7~\cite{ocmeIV2.7}, the Colorado Bureau of Investigation (CBI) Lab's internal validation of STRmix v2.5~\cite{cbiIV}, the Palm Beach Sheriff’s Office (PBSO) Lab's internal validation of STRmix v2.6.2~\cite{pbsoIV}, and the Maryland State Police's (MSP) internal validation of STRmix v2.9.1~\cite{mspIV}.\footnote{Because the creators of STRmix claim that the software is being used in 91 labs in the U.S.~\cite{strmix2025use}, we expect there to be somewhere on the order of 100 audit reports. It is important to note that at the time of our study, this database only contained five.} This dataset covers a range of software versions, state and city forensic labs of different sizes and geographies, and years (2016-2024).\footnote{The first U.S. forensic lab internal validation of STRmix was completed in 2014~\cite{butler2024history}.} In total, we analyzed over 280 pages of audit reports.


\subsection{Assessing Study Compliance}
\label{sec:methods-assessing-compliance}
The first author first assessed each audit report against relevant ASB 018 requirements to understand whether the audit could be interpreted as complying with the standard. Importantly, only the MSP lab's audit report (completed in 2024)~\cite{mspIV} claims compliance with ASB 018. Three audits were completed before the publication of ASB 018 in 2020 (OCME v2.4~\cite{ocmeIV}, CBI v2.5~\cite{cbiIV}, PBSO v2.6.2~\cite{pbsoIV}). While the OCME lab's v2.7~\cite{ocmeIV2.7} audit report was completed in 2021, it does not mention ASB 018. Thus, our focus is on interpreting whether labs' practices \textit{could be interpreted} as complying with the standard. 

Furthermore, of the three audit reports completed before 2020, none of the labs uploaded updated reports to the repository after 2020, despite ASB 018 advising labs to review their audits and supplementing where necessary~\cite{asb018}. This not only suggests that labs with pre-2020 audits believe their audits to be compliant with ASB 018, but also that ASB 018 did not change labs' existing practices. We further expand on this observation in Section~\ref{sec:disc-compatibility}.


\subsubsection{Identifying line-level requirements of ASB 018} 
\label{subsubsec:line-level-reqs}
Adopting an approach similar to \citet{lawrence2023bureaucratic}, we extracted all of the individual line-level requirements in ASB 018. We interpreted any action proceeded by ``shall'' as a line-level requirement, and focused only on requirements for internal validation studies. Next, we wrote out, as unambiguously as possible, how we interpreted each requirement. This step was crucial due to multiple ambiguities in ASB 018's language. In general, when we identified ambiguity in a requirement, we used the most permissive interpretation. For example, to assess compliance with ASB 018's requirement that labs ``shall evaluate both the appropriate sample types [...] and the number of samples within each type [and] shall base this evaluation on the intended application of the software'' (R4.1.6), we asked: ``Does any part of the report claim that the study's evaluated samples are appropriate for the intended application of the software?''  We provide additional details and the full list of our interpretations of each of the line-level requirements used in our study in Appendix~\ref{sec:appendix-b-interpretations}.


\subsubsection{Assessing study compliance with each of the requirements of ASB 018} 
For each audit, the first author read through the audit report, and resolved any confusions about report language and study design with the fourth author. Then, for each study, the first author identified a portion of the report that could be interpreted as fulfilling our interpretation of each of ASB 018's line-level requirements. 

Throughout this process, ASB 018's design also created challenges for comprehensively assessing compliance. For example, ASB 018 requires that samples studied ``represent (in terms of number of contributors, mixture ratios, and the total DNA [amount]), the range of actual casework samples intended for analysis with the system at the laboratory''~\cite{asb018}. While this requirement is not subject to the ambiguities discussed in Section~\ref{subsubsec:line-level-reqs}, it was difficult for us to comprehensively assess compliance with this requirement because none of the labs explicitly stated their lab's intended range of analysis. Since ASB 018 does not require that labs explicitly state the range of samples intended for analysis, we did not interpret reports' lack of specified ranges as non-compliance. Instead, following our permissive approach to assessing study compliance, we accepted any language in the report that suggests the lab chose samples representative of their intended range (e.g., report language concluding that the audit demonstrated the software was ``fit for its intended purpose in the lab'') as evidence of study compliance with this requirement. Table~\ref{appendix-table-compliance} in Appendix~\ref{appendix-b-compliance} documents sections of each audit report that we took to fulfill each of our interpretations of ASB 018's line-level requirements.


\subsection{Data Analysis}
\label{sec:methods-data-analysis}
Our study investigates the gaps between internal validation practices that ASB 018 envisions and the internal validation practices it enables, and how the design of ASB 018 enables these gaps. To do so, we first identified ASB 018's goals for compliant audits, i.e., what practices ASB 018 envisions or seeks to ensure about compliant studies. Next, we examined how labs have conducted audits in practice. Last, we returned to the standard and examined how its design enables the gaps between the envisioned and realized practices identified in the previous two steps.


We analyzed ASB 018 using five audit stages from \citet{ojewale2025towards}'s framework for stages of an audit: Audit Scope, Standards Identification, Performance Analysis, Audit Judgment, and Audit Communication. We used these stages to connect our study to existing literature and ensure our findings are relevant to other audit contexts, and we focused on this specific subset of stages given their relevance to the PGS context we are studying. For each audit stage, the first author identified and qualitatively coded relevant excerpts from the standard and the factsheet. We also coded these documents to investigate how ASB 018 envisions the purpose of audits. The first author then identified relationships between codes within each audit stage (and also for audit purpose) and grouped codes together into increasingly abstract themes through an iterative, bottom-up affinity diagramming process~\cite{beyer1999contextual}. The group then further refined these themes through discussion, focusing on themes that captured key meanings within each audit stage and assessing the fit of each theme with its corresponding audit stage. We used our themes for ASB 018's vision for audit purpose to complement our construction of themes for the standard's goals for each audit stage. For example, we drew on our interpretation of the standard's emphasis on the role of audits in ``establishing boundaries on software use'' in our interpretation of one of the standard's goals as ``expectations are falsifiable''.

Next, to explore potential gaps between laboratories' validation practices and ASB 018's goals, the first author identified audit excerpts relevant to each of the ASB 018 themes and inductively coded those excerpts. Instead of looking for specific practices, we grounded our interpretations in the practices observable in the report. For each audit stage, the team discussed the initial list of codes to ensure codes captured meaningful divergences from the standard's goals and then iteratively grouped codes into higher-level themes. 


Finally, we sought to understand how the design of the standard enables the gaps we identified in the previous step. For each ASB 018 theme, the first author identified relevant ASB 018 requirements. The first author analyzed these excerpts alongside the codes and high-level themes describing internal validation practices, and generated codes capturing specific characteristics of the requirements that enable observed patterns in validation study practices. The first author then iteratively grouped codes into themes. 


For example, under the Audit Scope stage, a theme representing ASB 018's goals was ``test broader sociotechnical system'' (corresponding codes included ``internal validation as providing assurance about system'' and ``acknowledging impact of user input''), a theme representing internal validation practices was ``testing software only, not including user'' (codes included ``using ground truth number of contributor value''), and a theme representing aspects of ASB 018's design that enabled the mismatch between goal and practice was ``isolated requirements'' (codes included ``addressing performance measurement and user input separately''). 


 
\subsection{Study Limitations and Future Work}
\label{sec:methods-limits-fw}
\subsubsection{Audit artifacts as proxies for audit practice}
We seek to understand current audit practices by examining how they are discussed and represented in publicly available audit reports. In doing so, our findings our based on accounts of practices. Because these reports are typically submitted in court, we expect them to accurately represent auditing practices. We hope that our findings can compel increased transparency about audit practices so future work can examine audit practices using a variety of sources and methods to investigate the extent to which audit reports authentically and fully capture audit practices.

\subsubsection{Representativeness}
We analyzed audit reports produced by labs of varying types (city, county, and state), geographies, and sizes (e.g. large and small labs), for a variety of versions of the STRmix probabilistic genotyping software. We observed many similarities in the language and structure of the audit reports we examined, and we attribute these similarities to labs' (i) reliance on SWGDAM validation guidelines, a widely used alternative source of guidance for internal validation studies, and (ii) use of a reporting template provided by STRmix developers. These observations suggest that our findings may extend to audits that also reference the SWGDAM guidelines and audits by labs that also use the STRmix software.


However, we caution against assuming that these findings are applicable to all validation studies. At present, while the creators of STRmix claim that it is being used in 91 labs in the US and 29 more internationally~\cite{strmix2025use}, few labs have made their audit reports publicly available, which limits the extent to which we can conclude how representative these reports are~\cite{butler2024dna}. We hope that our findings can support efforts underway that are advocating for greater transparency of audit reports (e.g., \cite{canellas2021defending, butler2024dna, abebe2022adversarial, jin2024beyond, krane2022using, matthews2019right, matthews2020trusted, kinshipproblem}) that would enable future work that can additionally leverage quantitative methods to explore the representativeness of our findings of the broader landscape of audit reports. 


\subsubsection{Our permissive approach to assessing study compliance}  
While we drew on our collective familiarity and expertise in PGS and validation studies, our view does not necessarily represent the view of those who might actually be tasked with assessing laboratory audits for compliance with ASB 018, such as organizations and auditors tasked with forensic laboratory oversight and accreditation. Our study takes a permissive approach, towards our goal of identifying what practices are \textit{possible} under ASB 018. This allows us to investigate the largest possible gaps between auditing practices and the standard’s desired practices, to raise awareness of what could possibly happen and to help safeguard against these possibilities. Future work can investigate how actors tasked with forensic laboratory oversight and accreditation interpret ASB 018's requirements and the extent to which their interpretations vary.


\subsubsection{Critically examining ASB 018's goals} 
Our study aims to assess current audit practices against ASB 018's goals for PGS audits, towards our goal of assessing ASB 018's effectiveness and how the design of the standard may undermine its effectiveness at ensuring compliant audits meet its goals. A critical area for future work to investigate, to complement our assessment of the ASB 018's effectiveness, is the degree to which ASB 018's own goals comprehensively assess PGS system effectiveness. 

\begin{figure*}[h!]
    \centering
    \includegraphics[width=\textwidth]{results-summarized.png}
    \caption{Summary of the gap between each ASB 018 goal (blue text starting with \textbf{G:}) and the corresponding compliant but unsatisfactory audit practices (red text starting with \textbf{P:}). We identify gaps between ASB 018's envisioned practices and real-world audit practices for each of five stages of the audit process, which we draw from ~\citet{ojewale2025towards}.}.
    \label{fig:results-summary}
    \Description{Flowchart with five stages flowing from left to right: Audit Scope, Standards Identification, Performance Analysis, Post-Audit Judgment, and Audit Communication. For each stage, we frame the audit stage as a question that auditors would ask at the audit stage, and then we summarize our findings on the gaps between ASB 018's goals and compliant but unsatisfactory audit practices. Each gap corresponds to a subsubsection of the Results. Question for Audit Scope is: "What is the audit target and what issues is the audit prioritizing?" The two Goal-Practice gaps under Audit Scope are (1) Goal: Sociotechnical system, Practice: Inadequate evaluation of software user; (2) Goal: Samples representative of casework samples, Practice: Convenience Datasets. Question for Standards Identification is: "What expectations is the audit target evaluated against?" The three Goal-Practice gaps under Standards Identification are (1) Goal: Define expectations for system performance, Practice: Few explicitly stated expectations; (2) Goal: Expectations match lab customers’ expectations, Practice: Labs and their customers may have conflicting expectations; and (3) Goal: Falsifiable expectations, Practice: Subjective, contradictory expectations. Question for Performance Analysis is: "What data about system performance does the audit produce?" The two Goal-Practice gaps under Performance Analysis are (1) Goal: Measure performance on a variety of sample types, Practice: Evaluate a narrow subset of possible sample types; and (2) Goal: Sufficient sample size to address variability, Practice: No discussion of sufficiency. Question for Post-Audit Judgment is: "How are evaluation results translated to judgments, consequences, or actions?" The two Goal-Practice Gaps under Post-Audit Judgment are: (1) Goal: Assess whether measured behavior is (un)acceptable, Practice: Frame system failures as expected and acceptable; and (2) Goal: Establish boundaries on software use based on unacceptable performance, Practice: No such boundaries. Question for Audit Communication is "How are audit data, methods, findings, and conclusions documented?" The one Goal-Practice gap under Audit Communication is (1) Goal: Supports scrutiny and reproducibility by third parties, Practice: Insufficient information.}
\end{figure*}

\subsection{Positionality}
All authors were either trained as researchers in or work in the United States. The first author is a researcher in algorithm auditing and human-computer interaction. The second author is a researcher trained in statistics, and the fifth author is a researcher in empirical machine learning who has also served as an AI and law advisor to numerous defense, civic, and government organizations. Collectively, these authors' background and familiarity with auditing literature informed their decision to analyze ASB 018 and audit reports through the lens of five stages of the audit process and shaped the patterns they identified in the standard's design and the auditing practices.  

The third and fourth authors are experts in forensic DNA profiling and PGS, and the third author is also a researcher trained in molecular biology and population genetics: they have written numerous articles on the subject in forensics literature and frequently testify as expert witnesses in criminal cases involving evidence produced by forensic lab use of PGS. Throughout their casework, they have examined dozens of laboratory PGS validation studies and have encountered a variety of laboratory uses of PGS. Our specific focus on ASB 018 is driven by their experiences observing laboratory references to ASB 018 in criminal cases and their perspectives on ASB 018 as a standard with qualities that meaningfully differentiate the standard from other existing PGS validation guidelines in forensic DNA. Their expertise and experiences with investigating real-world uses of PGS in criminal cases also focused our analysis of audit practices on practices they commonly observed that created meaningful risks to defendants subject to evidence produced using these software.

